% Ensamblado PARCIAL del artículo VII. No es una edición de entrega.
% Título y tres movimientos conservados en PROPUESTA_EDITORIAL.md.
% Núcleo común conservado literalmente y antecedentes copiados de V REV02.
% El cierre textual de inputs no sustituye la validación de todas sus pruebas.
% La amplitud material lleva seguimiento INTEGRACION-VII-AMPLITUD.
% No crear PDF final ni declarar autonomía completa a partir de este archivo.
% Artículo VII — primera redacción editorial, 10 de septiembre de 2026.
% Apertura incluida en el ensamblado REV05, aún sin PDF autónomo.
% Procedencia y condiciones: PLAN_RESIDENCIAS_DEMOSTRATIVAS.md.
% No contiene ni sustituye las demostraciones que debe reunir el cuerpo.

\begin{abstract}
Se desarrolla la relación entre memoria orientada, geometría gravitatoria
y evolución cosmológica desde la construcción APP--TRIT--TPK. El recorrido
demostrativo articula la ley de dilución \(a^{-6}\), la existencia y
unicidad del umbral de rebote en la realización homogénea especificada y la
persistencia local del rebote transversal bajo perturbaciones controladas.
La memoria conserva orientación, acarreo e historia durante el retorno de
fase; su transporte admite realizaciones operatorias y geométricas que
preservan esas distinciones. El paso hacia la dinámica gravitatoria reúne
la corriente de pantalla, el funcional de amplitud torsional, su
normalización y su compatibilidad con el refinamiento. En la acción mínima
de Cartan--Holst, la torsión queda determinada algebraicamente por la
corriente de espín y evoluciona con la geometría y los campos materiales.
Los balances cosmológicos conservan las condiciones de las fuentes y de
los acoplamientos en cada carta. El desarrollo se completa con la
información de horizonte, las operaciones de lectura del observador y la
respuesta relacional del estado conjunto.
\end{abstract}
\clearpage

\section{Introducción}
\label{vii:sec:introduccion}

La relación entre gravitación y memoria adquiere contenido matemático
cuando se especifican el estado que conserva la historia, la operación que
la transporta y la ecuación en la que interviene su realización geométrica.
En la construcción considerada, el retorno de fase conserva una
periodicidad observable mientras el registro incorpora orientación,
acarreo y profundidad. El problema gravitatorio se sitúa en el transporte
de esa información y en su contribución a las variables geométricas y
materiales de la evolución.

Tres resultados organizan la exposición: la dilución de la memoria, el
rebote de la realización homogénea y la persistencia local de su umbral
transversal. La primera relación se expresa mediante dos lectores de una
misma ruta,
\begin{equation}
 \frac{a_n}{a_{\mathrm{ref}}}=\varphi^{n/3},
 \qquad M_n=\varphi^{-2n},
 \qquad
 M_n=\left(\frac{a_n}{a_{\mathrm{ref}}}\right)^{-6}.
 \label{vii:eq:programa-dilucion}
\end{equation}
Aquí \(a_{\mathrm{ref}}>0\) fija la normalización del factor de escala.
La potencia expresa la relación entre el crecimiento de la realización
espacial y el carácter cuadrático del lector de memoria. La contribución
gravitatoria requiere, además, determinar la corriente y la amplitud que
acompañan a esa dependencia de escala. Ambas construcciones ocupan una
sección propia antes de introducirlas en las ecuaciones cosmológicas.

La Holografía Modular Triádica aporta el soporte y las operaciones de esta
genealogía. APP conserva las hojas aditiva y multiplicativa y sus datos
residuales; TRIT tipa los regímenes y las orientaciones; TPK transporta,
memoriza y prolonga el estado. Sus publicaciones internas transmiten a
las construcciones posteriores tanto los valores como las relaciones que
los determinan. La Mecánica Dimensional realiza esas estructuras mediante
secciones de acción, magnitudes geométricas y observables físicos.

El primer movimiento del artículo reúne construcción y geometría. La
descomposición del transporte distingue memoria simétrica y circulación
orientada. La representación del grupoide de rutas en geometría de Cartan
se desarrolla con composición, inversión, covariancia y refinamiento. La
acción, el parámetro de Barbero--Immirzi, la constante gravitatoria y la
unidad areal conservan sus antecedentes estructurales; su reutilización
transmite las condiciones que individualizan cada sección.

El segundo movimiento desarrolla amplitud, evolución y persistencia.
% INTEGRACION-VII-AMPLITUD: aplicar la corriente a la realización material
% y obtener su contribución a la densidad mediante la variación métrica.
% Seguimiento: PLAN_RESIDENCIAS_DEMOSTRATIVAS.md, sección 2;
% Ya incorporados: pantalla y mínimo (30), conexión (30d),
% diferencial general (30b/30c), campo y peso minimizantes con su
% diferencial (30e), inversas y respuesta algebraica de Cartan (31).
% No volver a registrar esas operaciones como pendientes.
La corriente de pantalla y la corriente variacional de materia se
presentan con el mapa que permite componerlas. La normalización tracial
fija la lectura de densidad correspondiente y se mantiene diferenciada de
una suma extensiva. La variación de la conexión determina la respuesta
torsional de la acción mínima. Al evolucionar la cotetrada y los campos
materiales, cambia también la torsión algebraicamente asociada a su
corriente de espín.

La formulación del rebote conserva las densidades positivas, el signo de
la contribución torsional y la desigualdad de exponentes de la realización
homogénea. Su persistencia local se demuestra mediante dos controles
separados: una cota uniforme de la perturbación conserva el cambio de
signo en los extremos de un intervalo; una cota uniforme de su derivada
conserva la monotonía y la transversalidad. Las realizaciones de los
selectores cosmológicos se enuncian con sus propios dominios, junto a los
invariantes y balances que producen.

La conservación del tensor residual se formula en la carta métrica
seleccionada, con acoplamientos constantes sobre ella y una fuente
material covariantemente conservada. Su descomposición en traza y parte
sin traza permite identificar el intercambio entre ambos sectores;
la traza constante caracteriza el caso de conservación separada.

El tercer movimiento estudia horizontes y respuesta relacional. El estado
reducido, los registros accesibles y la purificación permiten describir
qué información publica una lectura y qué información permanece en las
correlaciones del estado conjunto. La definición operacional del
observador comprende orientación y trivialización del transporte,
restricción a observables y proyección al subsistema. La respuesta
autoinercial se conecta con esas operaciones mediante la factorización
por la traza global y el defecto geométrico de proyección.

Los balances de horizonte conservan igualmente las convenciones de
área, energía, temperatura y estado que utilizan.

Cada desarrollo culmina con el objeto construido, la operación que lo
determina, su invariante y su función física dentro de la realización.
El orden de las pruebas sigue las dependencias de esos objetos; la
organización narrativa hace visible su relación con los tres resultados
principales desde el comienzo. Las conclusiones reúnen también la
operación del observador y la respuesta relacional, de modo que los
balances informacionales y la dinámica gravitatoria conserven su
conexión explícita.


\input{sections/nucleo.tex}
\input{sections/extension.tex}
\input{sections/generacion.tex}
\input{sections/registro_k.tex}
\input{sections/alpha.tex}
\input{sections/excepcional.tex}
\input{sections/v_accion_estructural.tex}
\input{sections/v_coordenadas_angulares.tex}
\input{sections/v_funcional_barbero.tex}
\input{sections/v_unidad_areal.tex}
\input{sections/v_area_informacion.tex}
\input{sections/v_transduccion_termica.tex}
\input{29_retorno_y_memoria_traslacional.tex}
\input{30_pantalla_y_respuesta.tex}
\input{30d_realizacion_geometrica.tex}
\input{30b_variacion_energia_memoria.tex}
\input{30c_composicion_corriente_conexion.tex}
\input{30e_corriente_extension_minima.tex}
\input{31_corriente_y_respuesta_cartan.tex}
\input{50_dilucion_y_rebote.tex}
\input{60_tensor_residual_y_balance.tex}
\input{70_horizontes_e_informacion.tex}
\input{80_observador_y_respuesta_relacional.tex}

